\documentclass[10pt,a4paper]{article}
\usepackage[T1]{fontenc}
\usepackage[utf8]{inputenc}
\usepackage[ngerman]{babel}
\usepackage[margin=2cm]{geometry}
\usepackage[hidelinks]{hyperref}
\setlength{\parindent}{0pt}
\setlength{\emergencystretch}{2em}
\renewcommand{\arraystretch}{1.15}
\begin{document}
\begin{center}
  {\Large\textbf{Webtechnologien -- Übungsblatt 1}}\\[4pt]
  Aufgabe 1: HTTP-Requests \quad -- \quad 11.10.2026
\end{center}

\subsection*{1a und 1b: AAU und ORF}
\textbf{Aufgabe:} HTTP-Requests von AAU und ORF untersuchen:
Anzahl, Datenmenge, Ladezeit, Statuscode, Webserver, Ressourcentypen,
langsamster Request und Cookies.

\textbf{Ergebnisse:}

\begin{center}
\begin{tabular}{|p{4.1cm}|p{5.5cm}|p{5.5cm}|}
\hline
\textbf{Messwert} & \textbf{aau.at} & \textbf{orf.at} \\ \hline
Requests & 206 & 85 \\ \hline
Übertragene Daten & 9,7 MB & 4700 kB \\ \hline
Ressourcengröße & 9,8 MB & 6284 kB \\ \hline
DOMContentLoaded & 1,76 s & 731 ms \\ \hline
Load & 2,22 s & 1,06 s \\ \hline
Finish & 9,56 s & 7,84 s \\ \hline
Erster Dokument-Request & GET, 200 OK & GET, 200 OK \\ \hline
Webserver & nginx & Apache \\ \hline
Ressourcentypen & HTML, CSS, JavaScript, Schriftarten, PNG, JPEG; eingebettete SVGs & HTML, CSS, JavaScript, Schriftarten, JPEG, Fetch/XHR \\ \hline
Langsamster Request & \nolinkurl{wpml-mod.min.js?ver=7.1.6}: 648 ms & \nolinkurl{semantic.js?tenantId=1655}: 185 ms (blockiert); erfolgreiches JPEG: 152 ms \\ \hline
Cookies & Vier, siehe unten & Keine beobachtet \\ \hline
\end{tabular}
\end{center}

\textbf{AAU-Cookies:}\\
\texttt{aviaCookieConsent},
\texttt{aviaPrivacyEssentialCookiesEnabled},\\
\texttt{aviaPrivacyMustOptInSetting},
\texttt{aviaPrivacyRefuseCookiesHideBar}.

{\small Die Request-Zahlen sind die angezeigten Gesamtwerte.
Im Mitschnitt waren auch Einträge von Browsererweiterungen sichtbar.}

\subsection*{1c: Google-Suche}
\textbf{Aufgabe:} Netzwerkaktivität beim Eintippen von
\glqq Klagenfurt\grqq{} in das Google-Suchfeld beobachten.

\textbf{Ergebnis:} Zehn XHR-Requests mit Status 200.
Der Parameter \texttt{q} enthielt
den jeweils eingegebenen Text. Die Antwort enthielt Suchvorschläge wie
\glqq klagenfurt wetter\grqq{} und \glqq klagenfurt einwohner\grqq.
Diese werden beim Tippen nachgeladen, ohne die ganze Seite neu zu laden.

\subsection*{1d: Cache}
\textbf{Aufgabe:} Cache, Löschen und Gründe für das Deaktivieren erklären.

\textbf{Ergebnis:} Der Browsercache speichert Dateien lokal zur
Wiederverwendung. Löschen ist über \glqq Cache leeren und vollständig
neu laden\grqq{} möglich; \textbf{Disable cache} deaktiviert ihn bei
geöffneten Entwicklertools. Beides hilft, nach Änderungen aktuelle Dateien
zu laden. Vergleichsmessung auf ORF:

\begin{center}
\begin{tabular}{|l|l|l|l|}
\hline
\textbf{Aufruf} & \textbf{Übertragen} & \textbf{Load} & \textbf{Finish} \\ \hline
Mit Cache & 3,3 MB & 1,22 s & 8,01 s \\ \hline
Cache deaktiviert & 5,5 MB (119 Requests) & -- & -- \\ \hline
Cache geleert und neu geladen & 5,4 MB & 10,51 s & 24,35 s \\ \hline
\end{tabular}
\end{center}

Mit Cache waren Einträge mit \texttt{memory cache} und
\texttt{disk cache} sichtbar.

\medskip
{\small \textbf{Hilfsmittel:} Codex für Erklärungen, Auswertung der eigenen Messwerte und Dokumentation.}
\end{document}
